\documentclass[14pt]{extarticle}

\usepackage{modules/mns}
\usepackage{mathspec}

\setmainfont[
	Path = f/,
	BoldFont=rb.ttf,
	ItalicFont=ri.ttf,
	BoldItalicFont=rbi.ttf
		]{r.ttf}
\setsansfont[
	Path = f/,
	BoldFont=rbb.ttf,
	ItalicFont=rli.ttf,
	BoldItalicFont=rbbi.ttf
		]{rl.ttf}
		
\setmathfont(Digits)[Path = f/]{rm.ttf}
\setmathfont(Latin)[Path = f/]{rmi.ttf}
\setmathfont(Greek)[Path = f/, Uppercase]{rm.ttf}
\setmathfont(Greek)[Path = f/, Lowercase]{rmi.ttf}

\setmonofont[Path = f/]{rm.ttf}

\newcolumntype{C}[1]{>{\centering\let\newline\\\arraybackslash\hspace{0pt}}m{#1}}

\begin{document}
\pagecolor{bcck}

\def\slide#1{\newpage \begin{center} \bfseries #1 \end{center}}
\def\c{\color{txxt}}
\def\dwt{\draw[very thick,color=txxt]}

\newpage
\ \vspace{-0.4cm}

\begin{center}
	\large Церемония награждения \\
	Олимпиады \\ [0.6cm]
	{\bf Математика НОН-СТОП \\ 2019 \\ [0.8cm]}
	{\normalsize 25 апреля \\ «Рекорд» \\}
\end{center}

\newpage
\ \vspace{-0.65cm}

\begin{center}
	Мероприятие проводится при поддержке \\ [1.25cm]
	\includegraphics[scale=1.45]{fpg_av/fpg}
\end{center}

\def\slovo#1#2{
\newpage
\begin{center}
	{\bf Слово предоставляется} \\ \vspace{1.1cm}
	\begin{tabular}{|C{7.6cm}|}
	\hline
		\makecell{\ \vspace{-0.3cm} \\ #1 \vspace{0.4cm}} \\
	\hline
		\makecell{\ \vspace{-0.3cm} \\ \bf\Large #2 \vspace{0.4cm}} \\
	\hline
	\end{tabular}
\end{center}
}

\slovo{Генеральный партнёр — \\ Кафедра математики \\ и информатики СПбАППО}{Е.\,Ю.~Лукичева}
\slovo{Председатель Жюри}{Д.\,Г.~Штукенберг}
\slovo{Площадка \\ {\bf ГБОУ СОШ 564} \\ Адмиралтейского района СПб}{Н.\,Л.~Корсакова}
\slovo{Площадка \\ {\bf ГБОУ СОШ 264} \\ Кировского района СПб}{Н.\,Б.~Симакова}
\slovo{Площадка \\ {\bf ГБОУ Гимназия 330} \\ Невского района СПб}{Ю.\,П.~Герке}
\slovo{Площадка \\ {\bf МАОУ Лицей 7} \\ города Бердска}{Т.\,А.~Смирнова}

\def\nagrpage#1#2{
\newpage
\begin{center}
	{\bf\large Награждение} \\ \vspace{1.1cm}
	\Large % \sffamily
	\begin{tabular}{|C{6.25cm}|}
	\hline
		\, #1
		\vphantom{$\int\limits_0^0$} \\
	\hline
		\, #2 класс
		\vphantom{$\int\limits_0^0$} \\
	\hline
	\end{tabular}
\end{center}
}

\def\nagrpagepro#1#2{
\newpage
\begin{center}
	{\bf\large Награждение} \\ \vspace{1.1cm}
	\Large % \sffamily
	\begin{tabular}{|C{6.25cm}|}
	\hline
		\, #1
		\vphantom{$\int\limits_0^0$} \\
	\hline
		\, #2 класс, про
		\vphantom{$\int\limits_0^0$} \\
	\hline
	\end{tabular}
\end{center}
}

\nagrpage{Похвал. отзывы}{4}
\nagrpage{Дипломы}{4}
\nagrpage{Похвал. отзывы}{5}
\nagrpage{Дипломы}{5}
\nagrpage{Похвал. отзывы}{6}
\nagrpage{Дипломы}{6}
\nagrpage{Похвал. отзывы}{7}
\nagrpage{Дипломы}{7}
\nagrpage{Похвал. отзывы}{8}
\nagrpage{Дипломы}{8}
\nagrpagepro{Похвал. отзывы}{7}
\nagrpagepro{Дипломы}{7}
\nagrpagepro{Похвал. отзывы}{8}
\nagrpagepro{Дипломы}{8}

\slide{Эта презентация}
\ \\ [0.3cm]

\begin{center}
	\url{http://bit.ly/mns-25apr}
\end{center}

\slide{«Ты школьник! — А может ты школьник?»}

Школьники всегда лгут, а студенты всегда говорят правду. Каждого попросили написать, сколько он видит студентов вокруг себя.

\ms Всегда ли можно понять, сколько студентов?

\slide{«Ты школьник! — А может ты школьник?»}

Нет, не всегда:

\ms 3, 3, 3, 3,\quad 5, 5, 5, 5, 5,\\ 1001, 1002, 1003, 1004.

\ms Студентов могло быть как 4, \\ так и 6.

\slide{Простые, но не простые-простые}

Для любого $n$ существует нату- \\ ральное число $N$, у которого ровно $n$ различных натуральных дели- \\ телей.

$$\c 2^{n-1}$$

Его делители — $2^0$, $2^1$, $\c\ldots$, $\c 2^{n-2}$, $\c 2^{n-1}$.

\slide{Розеттский камень}

Алфавит из $n$ букв: сколько букв \\ ни напиши одна поверх другой, \\ их все можно идентифицировать.

\ms Будем двигать по букве {\itshape засечку.}

\smallskip

\begin{center} \tikz{
	\foreach \x in {-3,...,3} {
		\dwt (\x cm, -0.4 cm) -- (\x cm, 0.4 cm);
		\dwt (\x cm - 0.1 cm, -0.4cm) -- (\x cm + 0.1 cm, -0.4cm);
		\dwt (\x cm - 0.1 cm, 0.4cm) -- (\x cm + 0.1 cm, 0.4cm);
		\dwt (\x cm - 0.1 cm, 0.1 * \x cm) -- (\x cm + 0.1 cm, 0.1 * \x cm);
	}
} \end{center}

\slide{Это кто? Это кот}

\smallskip

\begin{center} \tikz{
\begin{scope}[scale=0.46]
	\filldraw[thick,fill=txxt!30!bcck,draw=txxt] (-1.9,-2.8) rectangle (13.75,2.8);
	\filldraw[thick,fill=txxt!12!bcck,draw=txxt] (-0.5,-2.55) rectangle (13.5,2.55);
	\dwt (1,0) circle[radius=1cm]
		++(2.5,0) circle[radius=1.5cm]
		++(2,0) circle[radius=0.5cm]
		++ (2.5,0) circle[radius=2cm]
		++ (3,0) circle[radius=1cm]
		++ (1.5,0) circle[radius=0.5cm];
	\draw (-1.42,0) node[rotate=90]{\scriptsize Кольца для котов};
	\draw (-0.82,0) node[rotate=90]{\tiny\sffamily\itshape Арт.\,231421};
\end{scope}
} \end{center}

\ms Каждый сантиметр может либо \\ быть границей колец, либо \\ не быть. Поэтому ответ —
	$$2^{12}$$

\slide{Конференция анонимных геометров}

На каждой грани куба — по\\ $k$ любителей геометрии.

\ms Максимум — $6k$ человек. А минимум?

\slide{Конференция анонимных геометров}

Разместим людей в вершины, \\ они сразу будут занимать \\ по три места.

\ms Остальных будем сажать \\ на рёбра:
	$$\c 8 + \frac{6k-24}{2} = 3k-4.$$

\slide{Велопоход}

Проехали $1210$ км. Осталось втрое больше, чем проедет машина, в $4$ раза более быстрая, за время от сейчас до момента, когда останется столько же, сколько проехали.

\ms Пусть осталось $t$. \\ $\c t=12(t-1210)$. $\c t = 1320$.

\ms Расстояние — 2530 км.

\slide{Велопоход}

... Чтобы ответ в задаче составил 1400 км.	 \vspace{-0.2cm}

	$$\c t = 3c \cdot (t-A),\quad t = \frac{3cA}{3c-1}$$
	$$\c \text{Расстояние\ \ —\ \ } A \cdot \frac{6c-1}{3c-1}$$
	$$\c \text{Пример\ \ —\ \ } A=100,\ \ c = \frac{13}{36}$$

\slide{Я человек простой}

Дано произведение и среднее \\ арифметическое, можно ли \\ восстановить оценки?

\ms Да, можно: все, кроме $2$ и $4$, вы- \\ делим из разложения на простые.

\ms Заменами $\c 2,\ 2\ \rightarrow\ 4$ добьёмся \\ нужного среднего.

\slide{Оптимальск-на-Туломе}

1 м доски — 6000 руб. — \\ 2 табуретки

\ms 1 м доски + 4 м палок — \\ 10000 руб. — 4 табуретки

\ms Выгоднее пользоваться \\ вторым цехом, пока на это \\ есть деньги.

\slide{Оптимальск-на-Туломе}

Первому цеху нужно кратное \\ 0.5 количество метров доски.

\ms Значит, второму тоже. Он \\ сделает чётное число табуреток.

\slide{Оптимальск-на-Туломе}

\begin{center} \begin{tabular}{|l|}
\hline
	\makecell[l]{\ \vspace{-0.3cm} \\ 0.5 м доски — 3000 руб. — \\ 1 табуретка \vspace{0.4cm}} \\
\hline
	\makecell[l]{\ \vspace{-0.3cm} \\
		4.5 м доски + 18 м палок — \\
		45000 руб. — \\
		18 табуреток \vspace{0.4cm}
	} \\
\hline
\end{tabular} \end{center}

\slide{Меняем правила под себя}

$1$, $2$, $3$, $\c\ldots$, $\c a-1$, $\c\not{\!a}$, $\c a+1$, $\c\ldots$, $\c N-1$, \\ $N$ камней.

\ms В такой формулировке, если \\ $\c a \ne N$, сразу вытащим все камни. Иначе побеждает второй: \\ $\c x \rightarrow N-x$.

\slide{Меняем правила под себя}

$1$, $2$, $3$, $\c\ldots$, $\c a-1$, $\c\not{\!a}$, $\c a+1$, $\c\ldots$, $\c n-1$, \\ $n$ камней. Всего в кучке $N$.

\ms Если $\c a > \tfrac{n}{2}$, то второй выигрывает при $N$, кратном $a$, иначе выигрывает первый.

\slide{Меняем правила под себя}

Если $\c a \le \tfrac{n}{2}$, то второй \\ выигрывает при $N$ вида
	$$\c k \cdot (n+a+1),$$ \vspace{-0.6cm}
	$$\c k \cdot (n+a+1) + a.$$ \vspace{-0.6cm}
	
\ms Ответ получается методом \\ анализа позиций.

\def\tas{v_{_{\mathrm{TAS}}}} \def\gs{v_{_{\mathrm{GS}}}}

\slide{С вами говорит капитан}

Известно, что $\c \tas = 900\,\text{км/ч}$, \\ а скорость ветра $\c w =100\,\text{км/ч}$. \\
Куда дуть ветру, чтобы \\ $\c \gs = 900\,\text{км/ч}$?

\ms {\bf Не} в бок, потому что \\ Теорема Пифагора.

\slide{С вами говорит капитан}

$T$ — скорость в штиль, \\ $\c (a,b)$ — ветер. \vspace{-0.7cm}

\begin{center}\begin{tabular}{ll}
\def\TT{0.44 * 8.97216} \def\aa{0.44 * 0.999612} \def\bb{0.44 * -0.0278424}
\hspace{-0.6cm}\begin{tabular}{l}
\tikz{
	\draw [very thick, color=txxt, ->] (0,0) -- (0, \TT cm);
	\draw (0, 0.5 * \TT cm) node[left]{$T$};
	\draw [very thick, color=txxt, ->] (0, \TT cm) -- ++(\aa cm, 4*\bb cm);
	\draw (0.5 * \aa cm, \TT cm + 2*\bb cm) node[above]{$w$};
	\dwt (0,0) -- (\aa cm, \TT cm + 4*\bb cm);
	\draw (0.5 * \aa cm, 0.5*\TT cm + 2*\bb cm) node[right]{$\c\tas$};
}
\end{tabular}\hspace{-0.35cm}
&
\makecell[l]{
$\c
\begin{cases}
	T-b = 900 \\
	a^2 + b^2 = 100^2 \\
	(T-a)^2 + b^2 = 900^2
\end{cases}
$}
\end{tabular}\end{center}

\newpage
\ \\ [0.8cm]

\begin{center}
	\Large Спасибо за внимание!
\end{center}

\end{document}